% !TeX root = ../main.tex
% ---------------------------------------------------------------------------
% Drill and Blast Engineering
% Scope extension: not derived from the RSE3010 weekly material.
% Key sources: Persson, Holmberg & Lee; Langefors & Kihlstrom; Hustrulid;
% AS 2187.2 (explosives); Holmberg-Persson near-field vibration model.
% Related script: ../TBM_Vibration.py
% ---------------------------------------------------------------------------
\chapter{Drill and Blast Engineering}
\label{ch:drill-and-blast}


\section{Role and anatomy of blasting}
\label{sec:role-and-anatomy-of-blasting}

\subsection{The role of drill and blast}
\label{sec:dnb-role}
\scopenote{Where blasting remains the economic choice: hard rock, variable
geometry, short drives, caverns and shafts. Comparison against mechanical
excavation on cost, rate and flexibility.}

\subsection{Anatomy of a blast round}
\label{sec:blast-round}
\scopenote{Cut (burn, wedge, fan), stoping, contour and lifter holes; the
sequence by which a free face is created and progressively exploited.}

\section{Explosives and blast design}
\label{sec:explosives-and-blast-design}

\subsection{Explosives and initiation systems}
\label{sec:explosives}
\scopenote{Emulsions, ANFO, cartridged products; velocity of detonation,
energy and gas volume; electronic, non-electric and electric detonators, and
the precision that delay control now permits.}

\subsection{Blast design parameters}
\label{sec:blast-design}
\scopenote{Burden, spacing, hole diameter, charge concentration, specific
charge (powder factor) and drilling factor. Empirical design rules and their
calibration by trial rounds.}

\section{Perimeter control and excavation damage}
\label{sec:perimeter-control-and-excavation-damage}

\subsection{Perimeter control}
\label{sec:perimeter-control}
\scopenote{Smooth blasting, presplitting and cushion blasting; decoupled
charges; the trade-off between profile quality, cycle time and cost.}

\subsection{Overbreak and the blast damage zone}
\label{sec:overbreak}
\scopenote{Measuring overbreak, its cost in shotcrete and concrete, and the
extent of the damaged annulus. This is the physical basis of the disturbance
factor $\Dfac$ introduced in Chapter~\ref{ch:rock-mass-parameters}.}

\section{Vibration and environmental effects}
\label{sec:vibration-and-environmental-effects}

\subsection{Ground vibration prediction}
\label{sec:vibration-prediction}
\scopenote{Peak particle velocity, scaled distance laws, site-specific
regression, and the Holmberg--Persson near-field model for charges close to
the receiver.}

\subsection{Damage criteria for structures and utilities}
\label{sec:vibration-criteria}
\scopenote{Limits for buildings, heritage structures, buried services, fresh
concrete and sensitive equipment. Australian and international guidance, and
the difference between damage and disturbance thresholds.}

\subsection{Airblast, flyrock, fumes and ventilation}
\label{sec:blast-secondary-effects}
\scopenote{Overpressure limits, exclusion zones, post-blast fume clearance and
re-entry time as a constraint on cycle time.}

\section{Productivity and urban practice}
\label{sec:productivity-and-urban-practice}

\subsection{Cycle time and advance rate}
\label{sec:blast-cycle}
\scopenote{Drilling, charging, blasting, ventilation, mucking, scaling and
support as a cycle; where time is actually lost; realistic advance rates by
ground class.}

\subsection{Blasting in the urban environment}
\label{sec:urban-blasting}
\scopenote{Regulatory approvals, community notification, monitoring networks,
compliance reporting and dispute avoidance. Australian project practice.}
